\documentclass[11pt]{article}
\usepackage[T1]{fontenc}
\usepackage{amsmath,amssymb,amsthm}
\usepackage[margin=1in]{geometry}
\usepackage{microtype}
\usepackage{array}
\usepackage{needspace}
\usepackage[hidelinks]{hyperref}
\hypersetup{
  pdftitle={Quadratic cylinder energy and integrability of the Syracuse law},
  pdfauthor={Lucas Valbuena}
}
\newtheorem{theorem}{Theorem}[section]
\newtheorem{proposition}[theorem]{Proposition}
\newtheorem{lemma}[theorem]{Lemma}
\theoremstyle{remark}
\newtheorem*{remark}{Remark}
\newcommand{\Zthree}{\mathbb Z_3}
\newcommand{\E}{\mathbb E}
\newcommand{\Prob}{\mathbb P}
\newcommand{\one}{\mathbf 1}
\newcommand{\norm}[2]{\lVert #1\rVert_{#2}}
\DeclareMathOperator{\Law}{Law}
\title{Quadratic cylinder energy and integrability\\of the Syracuse law}
\author{Lucas Valbuena}
\date{8 October 2026}

\begin{document}
\maketitle

\begin{abstract}
Let \(\mu\) be the stationary probability measure on \(\Zthree\) for
the random affine maps \(x\mapsto(3x+1)/2^a\), where
\(\Prob(a=j)=2^{-j}\) for \(j\geq1\).  With respect to additive Haar
probability \(\lambda\), let \(\rho_k\) be the density obtained by
spreading each depth-\(k\) cylinder mass uniformly over that cylinder.
We prove
\[
  \int_{\Zthree}\rho_k^2\,d\lambda
  \leq24k+8k^2\leq32k^2\qquad(k\geq1).
\]
The proof combines offset injectivity, an elementary count of distinct
integer quotients, and normalized geometric second moments.  Together
with Tao's fine-scale mixing estimate, the bound gives a single density
\(\rho=d\mu/d\lambda\) in every \(L^p(\lambda)\), \(1\leq p<2\).
A support-overlap argument applied to the stationary equation shows
that \(\rho\notin L^2(\lambda)\).  These statements concern Haar
integrability of the random affine law.
\end{abstract}

\section{The stationary law and the results}

Write \(\lambda\) for additive Haar probability on \(\Zthree\), so
\(\lambda(x+3^k\Zthree)=3^{-k}\).  For \(a\geq1\), put
\[
  \Phi_a(x)=\frac{3x+1}{2^a}.
\]
If \(a_1,a_2,\ldots\) are independent with
\(\Prob(a_i=j)=2^{-j}\), the series
\begin{equation}\label{eq:random-series}
  X=\sum_{j=0}^{\infty}3^j2^{-(a_1+\cdots+a_{j+1})}
\end{equation}
converges in \(\Zthree\).  Its law \(\mu\) is stationary for the
maps \(\Phi_a\): separating the first exponent in the series gives
\(X\overset{d}{=}\Phi_{a_1}(X')\), where \(X'\) has law \(\mu\)
and is independent of \(a_1\).  This is the Syracuse law described in
\cite[Remark~1.13]{Tao}.  Its cylinder probabilities are
\[
  p_k(v)=\mu(v+3^k\Zthree),\qquad v\in\mathbb Z/3^k\mathbb Z,
\]
and we set
\begin{equation}\label{eq:rho-energy}
  \rho_k(x)=3^k p_k(x\bmod3^k),\qquad
  N_k=\int\rho_k^2\,d\lambda
      =3^k\sum_{v\bmod3^k}p_k(v)^2.
\end{equation}
In particular, \(\rho_0=1\) and \(N_0=1\).  All norms and integrals
below use the full additive Haar probability \(\lambda\).

\Needspace{6\baselineskip}
\begin{theorem}\label{thm:main}
The Syracuse law has the following properties.
\begin{enumerate}
\item For every integer \(k\geq1\),
\begin{equation}\label{eq:main-energy}
  N_k\leq24k+8k^2\leq32k^2.
\end{equation}
\item There is a nonnegative density \(\rho\), with
\(\mu=\rho\lambda\) and \(\int\rho\,d\lambda=1\), such that
\[
  \rho\in L^p(\lambda),\qquad
  \sup_{k\geq0}\norm{\rho_k}{p}<\infty
  \quad\text{for every }1\leq p<2.
\]
The same density \(\rho\) works for all these exponents.
\item This density does not belong to \(L^2(\lambda)\).
\end{enumerate}
\end{theorem}

The energy proof is independent of the mixing estimate.  The density
and its subcritical integrability use Tao's
Proposition~1.14~\cite{Tao}.  Offset injectivity, which enters the
energy argument, is the content of his Lemma~6.2; Section~6 of that
paper also uses separation of offsets in collision estimates.
We include the elementary decoding and moment arguments needed here.

These estimates are with respect to Haar measure.  Embedded ordinary
integers form a countable Haar-null set, so the estimates do not supply
an integer trace bound or a conclusion about individual Collatz orbits.

\section{Words, numerators, and fixed-total fibers}

For a word \(w=(a_1,\ldots,a_k)\) of positive integers, use
chronological composition:
\[
  \Phi_w=\Phi_{a_k}\circ\cdots\circ\Phi_{a_1},\qquad
  A_i=a_1+\cdots+a_i,\quad A_0=0,\quad A=A_k.
\]
Direct expansion gives
\begin{equation}\label{eq:word-formula}
  \Phi_w(x)=\frac{3^k x+B(w)}{2^A},\qquad
  B(w)=\sum_{i=0}^{k-1}3^{k-1-i}2^{A_i}.
\end{equation}
Define the positive real translation quantity and the residue
\begin{equation}\label{eq:translation-residue}
  Y(w)=\frac{B(w)}{3^k},\qquad
  r_k(w)=2^{-A}B(w)\pmod{3^k}.
\end{equation}
Here the inverse of \(2^A\) is taken modulo \(3^k\) in the
second expression.  The image \(\Phi_w(\Zthree)\) is the cylinder
\(r_k(w)+3^k\Zthree\).

The series~\eqref{eq:random-series}, reduced modulo \(3^k\), is the
endpoint of a length-\(k\) composition with the reverse ordering of
the exponents.  Reversal preserves their product distribution.
Consequently,
\begin{equation}\label{eq:fair-words}
  p_k(v)=\sum_{\substack{w\in\mathbb N_{>0}^k\\r_k(w)=v}}2^{-A(w)}.
\end{equation}
The sum is absolutely convergent since
\(\sum_w2^{-A(w)}=(\sum_{a\geq1}2^{-a})^k=1\).

\begin{lemma}[Fixed-total decoding]\label{lem:decoding}
For fixed \(k\geq1\) and \(A\), the numerator \(B(w)\)
determines a word \(w\) with \(A(w)=A\).  For a fixed residue
\(v\), the map
\[
  w\longmapsto q(w):=\left\lfloor\frac{B(w)}{3^k}\right\rfloor
\]
is therefore injective among words with \(A(w)=A\) and \(r_k(w)=v\).
Moreover, \(q(w)\) is a nonnegative integer and \(q(w)\leq Y(w)\).
\end{lemma}

\begin{proof}
For \(k=1\), the total fixes the only exponent.  For \(k\geq2\),
\eqref{eq:word-formula} gives
\[
  B(w)-3^{k-1}=2^{a_1}B(a_2,\ldots,a_k).
\]
The numerator on the right after removing \(2^{a_1}\) is odd:
its first summand is odd and its remaining summands are even.
Thus
\(a_1=v_2(B(w)-3^{k-1})\).  Division recovers the suffix
numerator, and induction with total \(A-a_1\) recovers the suffix.
This is the fixed-total numerator form of offset injectivity
\cite[Lemma~6.2]{Tao}.

On the fiber with total \(A\) and residue \(v\), every numerator
has the same remainder modulo \(3^k\), namely \(2^A v\).
Equal quotients therefore imply equal numerators and then equal words.
The last assertion follows from the definition of the floor.
\end{proof}

\begin{lemma}[Quotient count]\label{lem:quotient-count}
If \(q_1,\ldots,q_m\) are distinct nonnegative integers and
\(Y_j\geq q_j\), then
\[
  m^2\leq2\sum_{j=1}^m(1+Y_j).
\]
In particular, if \(m_A(v)\) counts any chosen subset of the
fixed-\((k,A,v)\) fiber, then
\begin{equation}\label{eq:fiber-count}
  m_A(v)^2\leq
  2\sum_{\substack{w\text{ in the chosen fiber}}}(1+Y(w)).
\end{equation}
\end{lemma}

\begin{proof}
Order the distinct integers increasingly.  The \(j\)-th is at
least \(j-1\), so
\[
  \sum_{j=1}^m(1+Y_j)\geq\sum_{j=1}^m j
  =\frac{m(m+1)}2\geq\frac{m^2}2.
\]
Apply Lemma~\ref{lem:decoding} to obtain~\eqref{eq:fiber-count}.
\end{proof}

\section{The energy estimate}

Fix \(k\geq1\).  Introduce the probability distribution on words
\begin{equation}\label{eq:biased-law}
  \pi_k(w)=3^k4^{-A(w)}
         =\prod_{j=1}^k(3\cdot4^{-a_j}).
\end{equation}
It is normalized because \(\sum_{a\geq1}3\cdot4^{-a}=1\).
This distribution will be used only to calculate a moment;
the cylinder probabilities in~\eqref{eq:fair-words} retain their fair
weights \(2^{-A(w)}\).

\begin{lemma}[Energy reduction]\label{lem:energy-reduction}
One has
\begin{equation}\label{eq:energy-moment}
  N_k\leq\frac4k\E_{\pi_k}\bigl[A^2(1+Y)\bigr].
\end{equation}
\end{lemma}

\begin{proof}
First restrict to words with \(A\leq M\), where \(M\geq k\).
For each residue \(v\), denote the corresponding partial sum in
\eqref{eq:fair-words} by \(p_{k,M}(v)\).  Every word has \(A\geq k\).
The elementary estimate
\[
  \frac1{A^2}\leq2\left(\frac1A-\frac1{A+1}\right)\qquad(A\geq1)
\]
implies \(\sum_{A=k}^{M}A^{-2}\leq2/k\).  Weighted
Cauchy--Schwarz and~\eqref{eq:fiber-count} now give
\begin{align*}
  p_{k,M}(v)^2
  &=\left(\sum_{A=k}^{M}2^{-A}m_A(v)\right)^2\\
  &\leq\frac2k\sum_{A=k}^{M}A^2 4^{-A}m_A(v)^2\\
  &\leq\frac4k\sum_{\substack{w:\,A(w)\leq M\\r_k(w)=v}}
                   A(w)^2 4^{-A(w)}(1+Y(w)).
\end{align*}
Sum over \(v\) and multiply by \(3^k\).  Each word has exactly
one residue, and hence
\[
  3^k\sum_v p_{k,M}(v)^2
  \leq\frac4k\sum_{A(w)\leq M}\pi_k(w)A(w)^2(1+Y(w)).
\]
Let \(M\to\infty\).  The left side tends to \(N_k\), and the
nonnegative sums on the right increase to the displayed expectation.
The next lemma shows that this expectation is finite.
\end{proof}

\begin{lemma}[Normalized second moments]\label{lem:moment}
For \(k\geq1\),
\begin{equation}\label{eq:moment-bound}
  \E_{\pi_k}\bigl[A^2(1+Y)\bigr]\leq6k^2+2k^3.
\end{equation}
\end{lemma}

\begin{proof}
The identity
\(\sum_{a\geq1}a^2z^a=z(1+z)/(1-z)^3\), for \(0<z<1\),
gives
\begin{equation}\label{eq:coordinate-moments}
  \sum_{a\geq1}a^2 2^{-a}=6,\qquad
  \sum_{a\geq1}a^2(3\cdot4^{-a})=\frac{20}{9}\leq6.
\end{equation}
For \(0\leq i\leq k\), put \(C_i=3^i2^{-A_i}\).  Formula
\eqref{eq:word-formula} yields
\begin{equation}\label{eq:translation-tilt}
  Y=\frac13\sum_{i=0}^{k-1}C_i^{-1}.
\end{equation}
Multiplying \(\pi_k\) by \(C_i^{-1}\) changes precisely the
first \(i\) coordinate weights:
\[
  (3\cdot4^{-a})\frac{2^a}{3}=2^{-a}.
\]
Therefore \(\pi_k(w)C_i(w)^{-1}\) is a normalized product
distribution, with fair coordinates before position \(i+1\) and
the weights \(3\cdot4^{-a}\) thereafter.  No normalization factor
is lost in this change of measure.

The pointwise inequality
\(A^2\leq k\sum_{j=1}^k a_j^2\), together with
\eqref{eq:coordinate-moments}, gives
\[
  \E_{\pi_k}[A^2C_i^{-1}]\leq6k^2
       \qquad(0\leq i\leq k).
\]
For \(i=0\), this bounds \(\E_{\pi_k}A^2\).
Summing~\eqref{eq:translation-tilt} gives
\[
  \E_{\pi_k}[A^2Y]
  =\frac13\sum_{i=0}^{k-1}\E_{\pi_k}[A^2C_i^{-1}]
  \leq2k^3.
\]
All the summands are nonnegative, so Tonelli's theorem justifies the
sums, and the bounds also establish their finiteness.
\end{proof}

Combining Lemmas~\ref{lem:energy-reduction} and~\ref{lem:moment}
proves
\[
  N_k\leq\frac4k(6k^2+2k^3)=24k+8k^2\leq32k^2,
\]
which is part~(1) of Theorem~\ref{thm:main}.

\section{Mixing and subcritical integrability}

The published input is the following consequence of
\cite[Proposition~1.14, equations~(1.26)--(1.27)]{Tao}:
for each \(A>0\) there is a finite constant \(C_A\) such that
\begin{equation}\label{eq:mixing}
  \norm{\rho_n-\rho_m}{1}\leq C_A m^{-A}
  \qquad(n\geq m\geq1).
\end{equation}
To check the normalization, consistency of the cylinder masses gives
\[
  \norm{\rho_n-\rho_m}{1}
  =\sum_{v\bmod3^n}
       \left|p_n(v)-3^{m-n}p_m(v\bmod3^m)\right|.
\]
This is exactly the oscillation in the cited proposition.

It follows that \(\rho_k\) is Cauchy in \(L^1(\lambda)\).
Let \(\rho\) be a nonnegative representative of its limit.
The limit has integral one.  For every cylinder \(D\) of depth
\(m\), consistency gives \(\int_D\rho_k\,d\lambda=\mu(D)\)
when \(k\geq m\).  Passing to the \(L^1\) limit proves the same
identity for \(\rho\).  Cylinders generate the Borel sigma algebra,
so the finite measures \(\rho\lambda\) and \(\mu\) are equal.

Fix \(1<p<2\), and put
\(\delta_k=\rho_{k+1}-\rho_k\) for \(k\geq1\).
By~\eqref{eq:mixing},
\begin{equation}\label{eq:increment-one}
  \norm{\delta_k}{1}\leq C_A k^{-A}.
\end{equation}
The elementary inequality \((u-v)^2\leq2(u^2+v^2)\) and the
energy bound give
\begin{equation}\label{eq:increment-two}
  \int\delta_k^2\,d\lambda
  \leq2(N_{k+1}+N_k)
  \leq64\bigl((k+1)^2+k^2\bigr)
  \leq320k^2.
\end{equation}
H\"older interpolation between the first and second moments yields
\begin{align}
  \norm{\delta_k}{p}
  &\leq\norm{\delta_k}{1}^{(2-p)/p}
            \left(\int\delta_k^2\,d\lambda\right)^{(p-1)/p}
       \notag\\
  &\leq C_A^{(2-p)/p}320^{(p-1)/p}
           k^{[-A(2-p)+2(p-1)]/p}.\label{eq:interpolation}
\end{align}
For completeness, the first inequality follows by writing
\(|u|^p=|u|^{2-p}(|u|^2)^{p-1}\), applying H\"older with
conjugate exponents \(1/(2-p)\) and \(1/(p-1)\), and taking
the \(p\)-th root.

Choose a positive integer
\[
  A>\frac{4p-2}{2-p}.
\]
Then the exponent in~\eqref{eq:interpolation} is at most \(-2\).
With \(D_p=C_A^{(2-p)/p}320^{(p-1)/p}\), finite telescoping and
Minkowski give, for every \(n\geq1\),
\begin{equation}\label{eq:uniform-lp}
  \norm{\rho_n}{p}
  \leq\norm{\rho_1}{p}+\sum_{k=1}^{n-1}\norm{\delta_k}{p}
  \leq\norm{\rho_1}{p}+2D_p.
\end{equation}
Here \(\rho_1\) is a bounded cylinder function, and
\(\sum_{k\geq1}k^{-2}\leq2\) follows from the same telescoping
estimate used in the energy proof.

The already established \(L^1\) convergence implies convergence in
measure to \(\rho\).  The Fatou property for convergence in measure
therefore gives
\[
  \norm{\rho}{p}\leq\liminf_{n\to\infty}\norm{\rho_n}{p}
  \leq\norm{\rho_1}{p}+2D_p.
\]
Indeed, one may choose a subsequence realizing the liminf, pass to a
further subsequence converging almost everywhere, and apply Fatou's
lemma to the \(p\)-th powers.  This step uses the same \(L^1\)
limit for each \(p\).  The case \(p=1\) was established above,
and \(\rho_0=1\) supplies the remaining finite level.
This proves part~(2) of Theorem~\ref{thm:main}.

\section{The stationary equation and the critical exponent}

Let
\[
  J(y)=\frac{3y+1}{2},\qquad
  D=J(\Zthree)=\{x\in\Zthree:x\equiv2\pmod3\}.
\]
Multiplication by \(2\) preserves \(\lambda\), while
\(J_*\lambda=3\lambda|_D\).  For a measurable function \(h\),
write its extension along this branch as
\[
  (Lh)(x)=
  \begin{cases}
    h((2x-1)/3),&x\in D,\\
    0,&x\notin D.
  \end{cases}
\]
Thus, whenever the integral is defined,
\begin{equation}\label{eq:lift-integral}
  \int Lh\,d\lambda=\frac13\int h\,d\lambda,
\end{equation}
and, for integrable \(h,g\),
\begin{equation}\label{eq:lift-l1}
  \norm{Lh-Lg}{1}=\frac13\norm{h-g}{1}.
\end{equation}

The geometric weights yield the exact finite-level recursion
\begin{equation}\label{eq:finite-recursion}
  \rho_{k+1}(x)
   =\frac12\rho_{k+1}(2x)+\frac32(L\rho_k)(x)
       \qquad(k\geq0).
\end{equation}
To see this, separate \(a=1\) from \(a\geq2\) in the outermost
exponent of the random series.  The first event has probability
\(1/2\) and applies \(J\) to the depth-\(k\) tail.  Conditional
on \(a\geq2\), the exponent \(a-1\) has the original geometric
law, and the output is half the corresponding depth-\((k+1)\)
variable.  The two changes of variables have Haar density factors
\(3\) and \(1\), respectively, which gives~\eqref{eq:finite-recursion}.

Passing to the \(L^1\) limit in~\eqref{eq:finite-recursion}, using
the invariance under doubling and~\eqref{eq:lift-l1}, proves
\begin{equation}\label{eq:stationary-density}
  \rho(x)=\frac12\rho(2x)+\frac32(L\rho)(x)
       \quad\text{for almost every }x.
\end{equation}
For an explicit estimate, the \(L^1\) norm of the difference between
the two sides is at most
\[
  \frac32\norm{\rho_{k+1}-\rho}{1}
  +\frac12\norm{\rho_k-\rho}{1},
\]
which tends to zero.

We use the following finite-measure observation.

\begin{lemma}[Positive support overlap]\label{lem:support}
Let \(f,a,b\) be nonnegative measurable functions on a finite
measure space \((\Omega,\nu)\).  Suppose \(f=a+b\) almost
everywhere,
\[
  \nu\{a>0\}=\nu\{f>0\},\qquad
  0<\int b\,d\nu<\infty,
\]
and \(ab\) is integrable.  Then \(\int ab\,d\nu>0\).
\end{lemma}

\begin{proof}
Since \(a\leq f\) almost everywhere, \(\{a>0\}\) is contained
in \(\{f>0\}\) up to a null set.  Equality of their finite
measures makes these sets equal up to a null set.  If
\(\int ab\,d\nu=0\), nonnegativity and integrability imply
\(ab=0\) almost everywhere.  But wherever \(b>0\), the equation
\(f=a+b\) gives \(f>0\), hence \(a>0\) outside a null set.
This forces \(b=0\) almost everywhere, contradicting its positive
integral.
\end{proof}

Suppose now that \(\rho\in L^2(\lambda)\), and put
\[
  a(x)=\frac12\rho(2x),\qquad b(x)=\frac32(L\rho)(x).
\]
Both functions belong to \(L^2\), their product is integrable,
and~\eqref{eq:stationary-density} gives \(\rho=a+b\) almost
everywhere.  Doubling preserves Haar measure, so
\[
  \lambda\{a>0\}=\lambda\{\rho>0\}.
\]
Equation~\eqref{eq:lift-integral} and \(\int\rho\,d\lambda=1\)
give \(\int b\,d\lambda=1/2\).  Lemma~\ref{lem:support}
therefore implies \(\int ab\,d\lambda>0\).

On the other hand, the same changes of variables yield
\[
  \int a^2\,d\lambda=\frac14\int\rho^2\,d\lambda,
  \qquad
  \int b^2\,d\lambda=\frac34\int\rho^2\,d\lambda.
\]
Squaring \(\rho=a+b\) and integrating would then give
\[
  \int\rho^2\,d\lambda
   =\int\rho^2\,d\lambda+2\int ab\,d\lambda,
\]
a contradiction.  This proves part~(3) of
Theorem~\ref{thm:main}.  Since \(\lambda\) is a probability measure,
\(L^p(\lambda)\subseteq L^2(\lambda)\) for every real \(p\geq2\).
Thus, among real exponents \(p\geq1\), the density belongs to
\(L^p(\lambda)\) exactly when \(p<2\).
The support argument uses no ergodicity
statement and no claim that \(\rho\) is positive almost everywhere
on a prescribed ambient set.

\clearpage
\section*{Formalization record}

The following table records the corresponding Lean endpoints.  The
energy proof uses the native reverse ordering of the word exponents;
an explicit reversal identity connects its numerator to
\eqref{eq:word-formula}.  Its moment proof uses the same geometric
second moments through a direct recursion on native words.

\begin{center}
\small
\begin{tabular}{@{}>{\raggedright\arraybackslash}p{0.23\textwidth}
  >{\raggedright\arraybackslash}p{0.52\textwidth}
  >{\raggedright\arraybackslash}p{0.19\textwidth}@{}}
\hline
Statement & Lean endpoint & Recorded status\\
\hline
Native quadratic energy &
\nolinkurl{FairEnergy.residue_energy_le_linear_quadratic} &
Strict compile and kernel check passed\\[5pt]
Haar quadratic energy &
\nolinkurl{canonicalPadicRho_sq_integral_le_thirtytwo_quadratic} &
Strict compile and kernel check passed\\[5pt]
One subcritical density &
\nolinkurl{canonicalSyracuseMeasure_has_subcritical_density} &
Strict compile and kernel check passed\\[5pt]
Uniform finite-level \(L^p\), \(1<p<2\), \(n\geq1\) &
\nolinkurl{exists_canonicalPadicRho_uniform_eLpNorm_bound} &
Strict compile and kernel check passed\\[5pt]
Actual stationary equation &
\nolinkurl{canonicalHaarDensity_stationary} &
Strict compile and kernel check passed\\[5pt]
Failure of \(L^2\) &
\nolinkurl{canonicalHaarDensity_not_memLp_two} &
Strict compile and kernel check passed\\[5pt]
One density with the exact \(L^p\) range &
\nolinkurl{canonicalSyracuseMeasure_has_sharp_lp_density} &
Strict compile and kernel check passed\\
\hline
\end{tabular}
\end{center}

Here \nolinkurl{FairEnergy} abbreviates
\nolinkurl{CollatzCylinderPacking.Arithmetic.FairEnergy}; the other
displayed endpoints are in \nolinkurl{Erdos1135.Tao}.
The uniform-bound endpoint is public; the extension to \(n=0\)
and \(p=1\) uses the finite norm at depth zero and the mass-one
identity at every depth.
The accompanying bundle freshly compiled all 362 internal source
modules and checked each with the official Lean kernel checker.
Its complete declaration audit was also compiled and kernel checked,
for 726 successful commands.  The audit covers 10,465 module/constant
rows, 10,445 distinct names and 736 encoded private rows;
only \nolinkurl{propext}, \nolinkurl{Classical.choice} and
\nolinkurl{Quot.sound} occur as axioms.  Complete external input
inventories agreed before and after the run.  The external Lean and
Mathlib libraries were fingerprinted and reused; their complete source
closures were not rebuilt.  Subsequent setup reuse and record verification
in fresh processes, including a relocated directory, also passed.
The bundle preserves the earlier failed strict-compilation attempt
and the 32 reviewed compatibility patches.  The fresh dependency-download
branch was not exercised.

\section*{Use of AI}
OpenAI Codex was used extensively for mathematical exploration,
development of the Lean proofs and verification scripts, automated
review, and preparation and revision of this manuscript.
The accompanying formal artifacts have the verification scope
recorded above.  Automated review does not constitute external peer
review or imply review or endorsement by the authors of the cited works.

\begin{thebibliography}{9}
\bibitem{Tao}
T.~Tao,
Almost all orbits of the Collatz map attain almost bounded values,
\emph{Forum of Mathematics, Pi} \textbf{10} (2022), e12.
\href{https://doi.org/10.1017/fmp.2022.8}
{doi:10.1017/fmp.2022.8}.
Author manuscript: \href{https://arxiv.org/abs/1909.03562v7}
{arXiv:1909.03562v7}.
\end{thebibliography}

\end{document}
